\phantomsection\label{finite-kernel-model}

In general, the spectral side of the trace formula tends to be the side which is more universal across examples. The idea is to 

WRITE THIS LIST BETTER AND HAVE THE REST OF THIS SECTION FOLLOW THIS LIST. 
\begin{enumerate}
    \item Given an orthonormal basis of eigenfunctions
    \item Compute a given inner product 
    \item Get the trace through a kernel thingy
\end{enumerate}

Let $(X,\mu)$ be a sigma-finite measure space and put $H=L^2(X,\mu;\mathbb C)$, with inner product and norm
\[
\begin{aligned}
\langle f,g\rangle_{L^2(X,\mu)}&=\int_X f(x)\overline{g(x)}\,d\mu(x),\\
\|f\|_{L^2(X,\mu)}&=\left(\int_X|f(x)|^2\,d\mu(x)\right)^{1/2}.
\end{aligned}
\]
In these subscripts, $L^2(X,\mu)$ denotes the complex space $H$. Our convention is linearity in the first argument. Fix an orthonormal family $u_1,\ldots,u_N$, choose measurable representatives, and let $V$ be its span. The case $N=0$ is allowed.

\begin{definition}[Finite kernel and associated operator]\label{finite-kernel-operator}
For $A=(a_{ij})\in M_N(\mathbb C)$, define
\[
K_A(x,y)=\sum_{i,j=1}^N a_{ij}u_i(x)\overline{u_j(y)},
\qquad
T_Af=\sum_{i,j=1}^N a_{ij}\langle f,u_j\rangle_{L^2(X,\mu)} u_i.
\]
\end{definition}

The function $K_A$ is an integral kernel for $T_A$:
\[
(T_Af)(x)=\int_X K_A(x,y)f(y)\,d\mu(y)
\]
for almost every $x$. Indeed, Cauchy--Schwarz gives
\[
\int_X|u_j(y)f(y)|\,d\mu(y)\leq\|u_j\|_{L^2(X,\mu)}\|f\|_{L^2(X,\mu)}=\|f\|_{L^2(X,\mu)},
\]
so every term is integrable and the finite sum may be integrated term by term. The same inequality and the triangle inequality show that
\[
\|T_Af\|_{L^2(X,\mu)}\leq\left(\sum_{i,j}|a_{ij}|\right)\|f\|_{L^2(X,\mu)}.
\]
Thus $T_A$ is bounded, has image in $V$, and vanishes on $V^\perp$. Since $T_Au_j=\sum_i a_{ij}u_i$, its restriction to $V$ has matrix $A$.

For a finite set $X=\{1,\ldots,N\}$ with counting measure and the point indicators as basis, $K_A(i,j)=a_{ij}$. The kernel formula then reduces to matrix multiplication. In general, its two variables play the roles of the output and input indices.

\begin{definition}[Finite-rank trace]\label{finite-kernel-trace-definition}
Define $\operatorname{tr}(T_A)=\operatorname{tr}(T_A|_V)$.
\end{definition}

\begin{proof}[Well-definedness]
This is independent of the basis by \hyperref[linear-algebra.trace-similarity]{similarity invariance}. It is also independent of the supporting space: if the same operator has such a description on $W$, compare both restrictions on $V+W$. Relative to
\[
V+W=V\oplus\bigl((V+W)\cap V^\perp\bigr),
\]
the matrix is that on $V$ with a zero block added. The analogous decomposition using $W$ proves equality of the two traces.
\end{proof}

\begin{proposition}[Finite-rank trace identity]\label{finite-kernel-trace-identity}
The specified diagonal $x\mapsto K_A(x,x)$ is integrable. If $\lambda_1,\ldots,\lambda_N$ are the eigenvalues of $A$, counted with algebraic multiplicity, then
\[
\boxed{\int_X K_A(x,x)\,d\mu(x)
=\operatorname{tr}(T_A)
=\sum_{j=1}^N\lambda_j.}
\]
\end{proposition}

\begin{proof}
For every $i,j$,
\[
\int_X|u_i(x)\overline{u_j(x)}|\,d\mu(x)
\leq\|u_i\|_{L^2(X,\mu)}\|u_j\|_{L^2(X,\mu)}=1.
\]
The diagonal is therefore a finite sum of integrable functions. Orthonormality gives
\[
\int_X K_A(x,x)\,d\mu(x)
=\sum_{i,j}a_{ij}\langle u_i,u_j\rangle_{L^2(X,\mu)}
=\sum_i a_{ii}
=\operatorname{tr}(A).
\]
The matrix of $T_A|_V$ is $A$, and the \hyperref[linear-algebra.trace-eigenvalue-sum]{trace--eigenvalue identity} supplies the last equality. For $N=0$, all expressions vanish.
\end{proof}

The eigenvalues in this statement belong to the finite-dimensional restriction. No diagonalizability assumption is needed. If $A$ is diagonal in the chosen orthonormal basis, the kernel has the particular form
\[
K_A(x,y)=\sum_j\lambda_j u_j(x)\overline{u_j(y)}.
\]
Integrating its diagonal replaces each factor $\int_X|u_j|^2\,d\mu$ by $1$. For $A=I_N$, the operator is orthogonal projection onto $V$, and its trace is $N$.

\begin{example}[A rank-one kernel]\label{finite-kernel-rank-one-example}

On $[0,1]$ with Lebesgue measure $dx$, the function $u(x)=\sqrt3\,x$ satisfies $\|u\|_{L^2([0,1],dx)}=1$. The matrix $(\alpha)$, for $\alpha\in\mathbb C$, gives
\[
K(x,y)=3\alpha xy,
\qquad
(Tf)(x)=3\alpha x\int_0^1 yf(y)\,dy.
\]
Here $Tu=\alpha u$, $T$ vanishes on $u^\perp$, and
\[
\int_0^1K(x,x)\,dx=3\alpha\int_0^1x^2\,dx=\alpha.
\]
For $\alpha=0$, the operator has rank zero; the identity is unchanged.
\end{example}

\begin{remark}[The choice of diagonal]\label{kernel-diagonal-representatives}

Changing the representatives of the finitely many $u_j$ changes the specified diagonal only on a null set, so its integral is well defined. An arbitrary representative of a kernel defined almost everywhere on $X\times X$ need not have this property.

For example, the zero function and the indicator of $\{(x,y):x=y\}$ agree almost everywhere on $[0,1]^2$ for two-dimensional Lebesgue measure. Both define the zero integral operator, but their diagonal integrals are respectively zero and one. Consequently, a measurable kernel's diagonal integral is not determined by its product-almost-everywhere class. The finite separated formula used above supplies a specified diagonal; any more general diagonal trace formula requires an appropriate representative and its own integrability justification.
\end{remark}
